\begin{lemma}
For every input \(\varepsilon>0\), there is a strictly smaller positive
parameter \(\varepsilon'< \varepsilon\), fixed constants
\(\varepsilon_1,\varepsilon_2\), and a threshold \(n_0=n_0(\varepsilon_2)\)
satisfying the hierarchy
\(\noderef{ParameterHierarchy}\).
\end{lemma}
\begin{proof}
Fix an input \(\varepsilon>0\).  Choose
\[
 \varepsilon'=\min\{\varepsilon/2,1/32\}.
\]
Then
\[
0<\varepsilon'< \varepsilon,\qquad
\varepsilon'\le\varepsilon,\qquad
\varepsilon'<\frac1{16};
\]
indeed \(\varepsilon/2>0\) and \(1/32>0\), while
\(\varepsilon'\le\varepsilon/2<\varepsilon\) and
\(\varepsilon'\le1/32<1/16\).  This gives the two strict upper bounds on the
first parameter required in
\(\noderef{ParameterHierarchy}\).

Next choose
\[
\varepsilon_1=
\frac12\min\left\{
\varepsilon',
\left(\frac{(\varepsilon')^3}{32}\right)^2,
\frac{(\varepsilon')^3}{40}
\right\}.
\]
This is positive and satisfies
\[
0<\varepsilon_1<\varepsilon'
\]
and
\[
8\sqrt{\varepsilon_1}+10\varepsilon_1\le \frac{(\varepsilon')^3}{2}.
\]
Indeed, the second entry in the minimum gives
\(8\sqrt{\varepsilon_1}\le(\varepsilon')^3/4\), and the third gives
\(10\varepsilon_1\le(\varepsilon')^3/4\).  Their sum is then at most
\((\varepsilon')^3/2\).  Then choose
\[
\varepsilon_2=\frac12\min\left\{
\varepsilon_1,\frac{\varepsilon_1}{30},1\right\}.
\]
This gives
\[
0<\varepsilon_2<\varepsilon_1
\]
and
\[
3\varepsilon_2\le \frac{\varepsilon_1}{10}.
\]
It also gives \(0\le\varepsilon_2\le1\), one of the eventual-comparison
conjuncts.

It remains only to choose the threshold.  Write \(\eta=\varepsilon'\).  For
each integer \(n\), put
\[
L=\log n,\qquad
k_R=(1+\eta)\sqrt{nL},\qquad
K=\noderef{TwoBiteNaturalIndependenceScale}(1+\eta,n),\qquad
k=(K:\mathbb R),
\]
\[
m_R=\noderef{TwoBiteReducedVertexCount}(n)=\frac{n}{L^2},\qquad
M=\noderef{TwoBiteNaturalReducedVertexCount}(n),\qquad
m=(M:\mathbb R),
\]
and
\[
p=\frac12\sqrt{\frac{L}{n}},\qquad
t_1=\frac{\sqrt{nL}}{\log L},\qquad
t_2=n^{1/4+\eta},\qquad
t_3=n^{2\eta}.
\]
All comparisons below are for an arbitrary \(n\) larger than the threshold
eventually chosen.  First choose \(N_{\mathrm{pos}}\) such that for all
\(n>N_{\mathrm{pos}}\),
\[
n>1,\qquad L>1,\qquad \log L>0,\qquad k_R>2,\qquad m_R>2,
\qquad p>0,\qquad t_1,t_2,t_3>0,
\]
and also
\[
\sqrt{L/n}\le1,\qquad \frac{k_R+1}{n}\le1,\qquad
\frac{\sqrt n}{2L^{3/2}}\ge1.
\]
These are finite requirements: \(L=\log n\to\infty\),
\(\log L\to\infty\), \(k_R=(1+\eta)\sqrt{nL}\to\infty\),
\(m_R=n/L^2\to\infty\), while
\[
\sqrt{L/n}\to0,\qquad
\frac{k_R+1}{n}
= (1+\eta)\sqrt{\frac Ln}+\frac1n\to0,\qquad
\frac{\sqrt n}{2L^{3/2}}\to\infty.
\]
The inequalities \(n>1\), \(L>1\), and \(\log L>0\) make
\(m_R\ge0\), \(k_R\ge0\), \(p\ge0\), and \(t_1,t_2,t_3>0\).  The displayed
bound \(\sqrt{L/n}\le1\) gives \(p\le1/2\).
The floor and ceiling inequalities from
\(\noderef{NaturalParameterApproximation}\), applied with
\(\kappa=1+\eta\) to the nonnegative real scales \(m_R\) and \(k_R\), give
exactly the rounding conjuncts required by
\(\noderef{ParameterHierarchyEventualComparisons}\):
\[
k_R\le k<k_R+1,\qquad
m\le m_R<m+1.
\]
The extra \(N_{\mathrm{pos}}\) condition \((k_R+1)/n\le1\) now gives
\(k<n\), hence \(K\le n\) as a natural number.
The auxiliary estimates used later follow from these same rounding bounds.
Since \(k_R>2\), the ceiling inequality gives
\[
k<k_R+1\le \frac32k_R<2k_R.
\]
Since \(m_R>2\), the floor inequality \(m_R<m+1\) gives
\[
m>m_R-1\ge \frac12m_R,
\]
and the other floor inequality gives \(m\le m_R\).  Thus we may also use
\[
k\le2k_R,\qquad m_R/2\le m\le m_R,\qquad 0<m.
\tag{R}
\]
Then (R) gives \(m\ge m_R/2\), and therefore
\[
2pm\ge 2p\cdot\frac{m_R}{2}=pm_R
=\frac12\sqrt{\frac Ln}\cdot\frac{n}{L^2}
=\frac{\sqrt n}{2L^{3/2}}\ge1. \tag{P}
\]
This proves the lower-bound conjunct \(1\le2pm\) in
\(\noderef{ParameterHierarchyEventualComparisons}\).
Since \(K\le n\), the elementary finite-set bound gives
\[
\binom{n}{K}
=\texttt{Nat.choose}(n,K)
\le n^K=\exp(kL).
\tag{B}
\]
Finally, by \(\noderef{RealChooseTwo}\), if
\[
C=\binom{k}{2}_{\mathbb R}=\frac{k(k-1)}2,
\]
then \(k>2\) gives
\[
\frac{k^2}{4}\le C\le \frac{k^2}{2}.
\tag{C}
\]

Every remaining threshold below is finite by the elementary growth rule
\[
\frac{(\log n)^a(\log\log n)^b}{n^c}\to0
\quad(c>0)
\]
and its reciprocal form \(n^c/((\log n)^a(\log\log n)^b)\to\infty\).  The
proof nevertheless records the concrete ratio used for each conjunct, so the
threshold for that conjunct is obtained by requiring that displayed ratio to
be at most the indicated fixed constant.

We now choose one finite threshold for each remaining conjunct in
\(\noderef{ParameterHierarchyEventualComparisons}\).  For
\[
n^{4\eta}k\le \frac{\varepsilon_1}{2}k^2,
\]
divide by \(k^2>0\).  By \(k\ge k_R\), the ratio is at most
\[
\frac{n^{4\eta}}{k_R}
=\frac{1}{1+\eta}\frac{n^{-1/2+4\eta}}{\sqrt L}\to0,
\]
because \(\eta<1/16\).  Choose \(N_1\) making it at most
\(\varepsilon_1/2\).

For
\[
2kL^2\le \frac{\varepsilon_1}{8}k^2,
\]
the ratio is \(2L^2/k\le2L^2/k_R\), and
\[
\frac{2L^2}{k_R}=\frac{2}{1+\eta}\frac{L^{3/2}}{\sqrt n}\to0.
\]
Choose \(N_2\) making the ratio at most \(\varepsilon_1/8\).

For the three-term \(\noderef{RealChooseTwo}\) inequality, since \(L>1\),
\[
\frac{1}{\sqrt L}-\frac{1}{2L}\ge \frac{1}{2\sqrt L}.
\]
Thus
\[
\frac{C}{2L}+2kL^2+\frac{C}{\sqrt L}\le \frac{2C}{\sqrt L}
\]
follows from \(2kL^2\le C/(2\sqrt L)\).  Using \(C\ge k^2/4\), it is enough to
require \(16L^{5/2}\le k\).  Since \(k\ge k_R\) and
\[
\frac{16L^{5/2}}{k_R}
=\frac{16}{1+\eta}\frac{L^2}{\sqrt n}\to0,
\]
choose \(N_3\) making this true.

For the McDiarmid/union-bound inequality
\[
\exp\!\left(-\frac{C^2}{Lm n^{8\eta}}\right)\binom{n}{K}\le \frac1n,
\]
use \(C\ge k^2/4\), \(k\ge k_R\), and \(m\le m_R=n/L^2\).  Then the exponent
magnitude is at least
\[
\frac{k^4}{16Lm n^{8\eta}}
\ge
\frac{k_R^4}{16L(n/L^2)n^{8\eta}}
=\frac{(1+\eta)^4}{16}n^{1-8\eta}L^3.
\]
By (B), the logarithm of \(\texttt{Nat.choose}(n,K)\) is at most \(kL\).
The target margin in the exponent is \(2kL+L\).  Using \(k\le2k_R\), and
increasing the present threshold so that \(L\le(1+\eta)\sqrt n\,L^{3/2}\), we
have
\[
2kL+L\le4k_RL+L\le5(1+\eta)\sqrt n\,L^{3/2}.
\]
The explicit ratio of the exponent lower bound to this required margin is
\[
\frac{\frac{(1+\eta)^4}{16}n^{1-8\eta}L^3}
     {5(1+\eta)\sqrt n\,L^{3/2}}
=\frac{(1+\eta)^3}{80}n^{1/2-8\eta}L^{3/2}\to\infty
\]
because \(\eta<1/16\).  Choose \(N_4\) so that the exponent magnitude is at
least \(2kL+L\).  Then the product is at most
\(\exp(-(2kL+L))\exp(kL)\le \exp(-L)=1/n\).

For
\[
\frac{2}{\sqrt L}C\le \frac{\varepsilon_1}{2}k^2,
\]
use \(C\le k^2/2\).  The left side is at most \(k^2/\sqrt L\), so it is enough
to have \(1/\sqrt L\le\varepsilon_1/2\).  Choose \(N_5\) with this property.

For
\[
4\log k\le \frac{\eta^3}{2}pk^2,
\]
we have \(k\ge k_R\), \(k\le2k_R\), and
\[
pk_R^2=\frac{(1+\eta)^2}{2}\sqrt n\,L^{3/2}.
\]
Also \(\log k\le\log(2(1+\eta)\sqrt{nL})\) by \(k\le2k_R\).  Hence
\[
\frac{4\log k}{pk^2}
\le
\frac{4\log(2(1+\eta)\sqrt{nL})}{p k_R^2}\to0.
\]
Choose \(N_6\) making this ratio at most \(\eta^3/2\).

For
\[
t_1k\le\varepsilon_1k^2,
\]
divide by \(k^2\) and use \(k\ge k_R\):
\[
\frac{t_1}{k}\le\frac{t_1}{k_R}
=\frac{1}{(1+\eta)\log L}\to0.
\]
Choose \(N_7\) making this at most \(\varepsilon_1\).

For the huge-family overlap comparison, set \(H=100L^3\).  First record a
common bound for all the huge-projection finite comparisons.  Include in
\(N_{7a}\) the finite requirement \(\log L\ge1\).  Then by \(k\le2k_R\),
\(\eta<1\), and \(t_1=\sqrt{nL}/\log L\),
\[
\frac{2k}{t_1}
\le \frac{4(1+\eta)\sqrt{nL}}{\sqrt{nL}/\log L}
=4(1+\eta)\log L
\le8\log L.
\]
Therefore
\[
\frac{2k}{t_1}+1\le 9\log L,
\qquad
\binom{2k/t_1+1}{2}_{\mathbb R}
\le \frac{(9\log L)^2}{2}
=\frac{81}{2}(\log L)^2.
\tag{H}
\]
It is enough for the overlap conjunct to prove
\[
\binom{2k/t_1+1}{2}_{\mathbb R}H\le\frac12 k.
\]
By (H), the left side is at most \(4050L^3(\log L)^2\).  Since
\(k\ge k_R=(1+\eta)\sqrt{nL}\), the ratio of this explicit upper bound to
\(\frac12 k_R\) is
\[
\frac{4050L^3(\log L)^2}{\frac12(1+\eta)\sqrt{nL}}
=\frac{8100}{1+\eta}\frac{L^{5/2}(\log L)^2}{\sqrt n}\to0.
\]
Choose \(N_{7a}\) making this ratio at most \(1\).  Then the displayed
overlap inequality holds for all \(n>N_{7a}\).

For the huge same-color projection-size comparison,
\[
\frac{2k}{t_1}\,2pm\le \frac{\varepsilon_1}{10}k,
\]
divide by \(k>0\).  Using \(m\le m_R=n/L^2\) and
\(p=\frac12\sqrt{L/n}\), the ratio of the left side to \(k\) is
\[
\frac{4pm}{t_1}
\le
\frac{4\cdot\frac12\sqrt{L/n}\cdot n/L^2}{\sqrt{nL}/\log L}
=\frac{2\log L}{L^2},
\]
which tends to \(0\).  Choose \(N_{7b}\) making this ratio at most
\(\varepsilon_1/10\).

For the huge same-color projected-pair comparison,
\[
\frac{2k}{t_1}\binom{2pm}{2}_{\mathbb R}\le \varepsilon_1k^2,
\]
use \(m\le m_R=n/L^2\) and \(p=\frac12\sqrt{L/n}\).  Then
\[
2pm\le \sqrt{\frac{n}{L^3}},
\qquad
\binom{2pm}{2}_{\mathbb R}\le \frac{n}{2L^3}
\]
for all \(n>N_{\mathrm{pos}}\), because \(2pm\ge0\) and
\(\binom{x}{2}_{\mathbb R}\le x^2/2\) for \(x\ge0\).  Together with
\(2k/t_1\le8\log L\), the left side is at most
\[
8\log L\cdot \frac{n}{2L^3}=\frac{4n\log L}{L^3}.
\]
Since \(k^2\ge k_R^2=(1+\eta)^2nL\), the ratio of this upper bound to
\(k_R^2\) is
\[
\frac{4n\log L/L^3}{(1+\eta)^2nL}
=\frac{4}{(1+\eta)^2}\frac{\log L}{L^4}\to0.
\]
Choose \(N_{7e}\) making this ratio at most \(\varepsilon_1\).

For the projected-overlap size correction, (H) gives the explicit upper bound
\[
\binom{2k/t_1+1}{2}_{\mathbb R}200L^3
\le 8100L^3(\log L)^2.
\]
The ratio of this upper bound to \((\varepsilon_1/10)k_R\) is
\[
\frac{8100L^3(\log L)^2}
     {(\varepsilon_1/10)(1+\eta)\sqrt{nL}}
=\frac{81000}{\varepsilon_1(1+\eta)}
  \frac{L^{5/2}(\log L)^2}{\sqrt n}\to0.
\]
Choose \(N_{7c}\) making this ratio at most \(1\).  Then
\[
\binom{2k/t_1+1}{2}_{\mathbb R}200L^3
  \le \frac{\varepsilon_1}{10}k
\]
for all \(n>N_{7c}\).
For the projected-pair multiplicity correction,
\[
\binom{2k/t_1+1}{2}_{\mathbb R}\binom{200L^3}{2}_{\mathbb R}
  \le \frac{\varepsilon_1}{10}k^2,
\]
use (H) and, for \(200L^3\ge0\),
\[
\binom{200L^3}{2}_{\mathbb R}
\le \frac{(200L^3)^2}{2}=20000L^6.
\]
Thus the left side is at most \(810000L^6(\log L)^2\).  The ratio of this
bound to \((\varepsilon_1/10)k_R^2\) is
\[
\frac{810000L^6(\log L)^2}
     {(\varepsilon_1/10)(1+\eta)^2nL}
=\frac{8100000}{\varepsilon_1(1+\eta)^2}
  \frac{L^5(\log L)^2}{n}\to0.
\]
Choose \(N_{7d}\) making this ratio at most \(1\), which proves the displayed
inequality.

We also impose the relative huge-error comparisons used when the \(o(k)\)
losses in the huge opposite-color argument are converted into
deficit-relative errors.  First,
\[
2pm\le\frac{\varepsilon_2}{100}t_1
\]
follows because, using \(m\le m_R=n/L^2\),
\[
\frac{2pm}{t_1}
\le
\frac{\sqrt{L/n}\cdot n/L^2}{\sqrt{nL}/\log L}
=\frac{\log L}{L^2}\to0.
\]
Choose \(N_{7f}\) making this ratio at most \(\varepsilon_2/100\).

Next, by \(2k/t_1+1\le9\log L\),
\[
\frac{(2k/t_1+1)100L^3}{t_1}
\le
\frac{900L^3\log L}{\sqrt{nL}/\log L}
=900\frac{L^{5/2}(\log L)^2}{\sqrt n}\to0,
\]
and the same estimate with \(100L^3\) replaced by \(200L^3\) differs only by a
factor \(2\).  Choose \(N_{7g}\) making both ratios at most
\(\varepsilon_2/100\).

Finally,
\[
\frac{(2k/t_1+1)400L^5}{t_1}
\le
\frac{3600L^5\log L}{\sqrt{nL}/\log L}
=3600\frac{L^{9/2}(\log L)^2}{\sqrt n}\to0.
\]
Choose \(N_{7h}\) making this ratio at most \(\varepsilon_2/100\).
These three threshold choices prove the four new relative comparison conjuncts.

Finally,
\[
3\varepsilon_2\le\frac{\varepsilon_1}{10}
\]
was one of the two explicit requirements in the initial choice of
\(\varepsilon_2\).  Also \(0<\varepsilon_2<\varepsilon_1<\eta<1\), so
\(0\le\varepsilon_2\le1\).  These two \(\varepsilon_2\)-conjuncts therefore
hold for every \(n\), independently of the threshold.

For
\[
t_2n^{1/4-\eta}\le\varepsilon_1k,
\]
the left side is \(n^{1/2}\), and
\[
\frac{n^{1/2}}{k}\le \frac{n^{1/2}}{k_R}
=\frac{1}{(1+\eta)\sqrt L}\to0.
\]
Choose \(N_8\) making this at most \(\varepsilon_1\).

For the cutoff comparison
\[
\noderef{TwoBiteLargeCutoff}(\eta,n)<\noderef{TwoBiteHugeCutoff}(n),
\]
that is, \(t_2<t_1\), compute
\[
\frac{t_2}{t_1}
=\frac{n^{1/4+\eta}}{\sqrt{nL}/\log L}
=n^{-1/4+\eta}L^{-1/2}\log L.
\]
Since \(\eta<1/16<1/4\), this ratio tends to \(0\).  Choose \(N_9\) so that
the ratio is \(<1\) for all \(n>N_9\).  Then
\[
\noderef{TwoBiteLargeCutoff}(\eta,n)=t_2
<t_1=\noderef{TwoBiteHugeCutoff}(n).
\]

For the two-term cutoff comparison
\[
(1+\varepsilon_2)L^2\frac{t_2}{L}
  +L(1+\varepsilon_2)pm<t_1,
\]
we do not use only the simpler inequality \(t_2<t_1\).  Instead, use
\(0\le\varepsilon_2\le1\), so \(1+\varepsilon_2\le2\), and estimate the two
ratios to \(t_1\) separately.  The first ratio is
\[
\frac{(1+\varepsilon_2)Lt_2}{t_1}
\le
2\frac{L\,n^{1/4+\eta}}{\sqrt{nL}/\log L}
=2n^{-1/4+\eta}L^{1/2}\log L\to0,
\]
because \(\eta<1/16<1/4\).  The second ratio, using \(m\le m_R=n/L^2\), is
\[
\frac{(1+\varepsilon_2)Lpm}{t_1}
\le
2\frac{L\cdot\frac12\sqrt{L/n}\cdot n/L^2}
        {\sqrt{nL}/\log L}
=\frac{\log L}{L}\to0.
\]
Choose \(N_{9a}\) making each displayed ratio strictly less than \(1/2\).
Then the two summands have total size strictly less than \(t_1\), which is
exactly this conjunct of
\(\noderef{ParameterHierarchyEventualComparisons}\).

For the monotonicity-size comparison used in the large-huge count,
\[
100\left(n^{1/4}+\frac12\right)L^3
\le \noderef{TwoBiteLargeCutoff}(\eta,n)=t_2,
\]
divide by \(t_2=n^{1/4+\eta}\):
\[
\frac{100(n^{1/4}+\frac12)L^3}{t_2}
=100\left(1+\frac{1}{2n^{1/4}}\right)L^3n^{-\eta}\to0.
\]
Choose \(N_{10}\) making this ratio at most \(1\).  Then for all
\(n>N_{10}\),
\[
100\left(n^{1/4}+\frac12\right)L^3\le t_2
=\noderef{TwoBiteLargeCutoff}(\eta,n).
\]

For the large-huge count comparison,
\[
k<n^{1/4}t_2-\binom{n^{1/4}}{2}_{\mathbb R}100L^3,
\]
the main term is \(n^{1/4}t_2=n^{1/2+\eta}\).  Since
\(\binom{x}{2}_{\mathbb R}\le x^2/2\) for \(x=n^{1/4}\ge1\), the subtracted
term is at most \(50n^{1/2}L^3\).  The two ratios
\[
\frac{50n^{1/2}L^3}{n^{1/2+\eta}}=50L^3n^{-\eta},
\qquad
\frac{k}{n^{1/2+\eta}}\le 2(1+\eta)\sqrt L\,n^{-\eta}
\]
both tend to \(0\).  Choose \(N_{11}\) so that the first ratio is at most \(1/4\)
and the second is at most \(1/4\).  Then the right side is at least
\((3/4)n^{1/2+\eta}\), while \(k\le(1/4)n^{1/2+\eta}\), proving the strict
inequality.

For
\[
t_3^2k\le \frac{\varepsilon_1}{2}k^2,
\]
we have \(t_3^2=n^{4\eta}\), so the same estimate as for \(N_1\) applies.
Choose \(N_{12}\) making \(t_3^2/k\le\varepsilon_1/2\).

For
\[
k\left(\frac{2k}{L}+n^{1/4}\right)\le\varepsilon_1k^2,
\]
divide by \(k^2\):
\[
\frac{2}{L}+\frac{n^{1/4}}{k}
\le
\frac2L+\frac{n^{1/4}}{k_R}
=\frac2L+\frac{1}{(1+\eta)n^{1/4}\sqrt L}\to0.
\]
Choose \(N_{13}\) making this at most \(\varepsilon_1\).

For
\[
8\sqrt{\varepsilon_1}\,pk^2+10\varepsilon_1pk^2+4\log k
  \le \eta^3pk^2,
\]
the parameter choice gives
\[
8\sqrt{\varepsilon_1}+10\varepsilon_1\le\frac{\eta^3}{2},
\]
and the threshold \(N_6\) gives \(4\log k\le(\eta^3/2)pk^2\).  Therefore this
conjunct holds for all \(n>N_6\); set \(N_{14}=N_6\).

For the medium exponential comparison,
\[
\exp\!\left(kL+2kn^{1/4-3\eta}L-n^{3/4-2\eta}\right)\le\varepsilon_1,
\]
let
\[
E(n)=kL+2kn^{1/4-3\eta}L-n^{3/4-2\eta}.
\]
Using \(k\le2k_R\), the first positive term has ratio at most
\[
\frac{2k_RL}{n^{3/4-2\eta}}
=2(1+\eta)L^{3/2}n^{-1/4+2\eta}\to0,
\]
and the second has ratio at most
\[
\frac{4k_Rn^{1/4-3\eta}L}{n^{3/4-2\eta}}
=4(1+\eta)L^{3/2}n^{-\eta}\to0.
\]
Choose \(N_{15}\) so that the sum of these ratios is at most \(1/2\), giving
\[
E(n)\le-\frac12 n^{3/4-2\eta}.
\]
Increase \(N_{15}\), if necessary, so that
\(\exp(-\frac12 n^{3/4-2\eta})\le\varepsilon_1\).

For the final exponential comparison,
\[
\exp\!\left(-\frac{\eta}{4}kL\right)\le\varepsilon_1,
\]
we have \(kL\ge k_RL=(1+\eta)\sqrt n\,L^{3/2}\to\infty\).  Choose \(N_{16}\)
so that \((\eta/4)kL\ge\log(1/\varepsilon_1)\).

There are finitely many thresholds
\[
N_{\mathrm{pos}},N_1,\ldots,N_{16},
N_{7a},N_{7b},N_{7c},N_{7d},N_{7e},N_{7f},N_{7g},N_{7h},N_{9a}.
\]
Choose \(n_0\) to be an integer at least
\[
\max\{\varepsilon_2^{-2},N_{\mathrm{pos}},N_1,\ldots,N_{16},
N_{7a},N_{7b},N_{7c},N_{7d},N_{7e},N_{7f},N_{7g},N_{7h},N_{9a}\}.
\]
Then \(n_0\) satisfies the lower bound in
\(\noderef{ParameterHierarchy}\), and for every \(n>n_0\) every conjunct of
\(\noderef{ParameterHierarchyEventualComparisons}(\eta,\varepsilon_1,
\varepsilon_2,n_0)\) holds by the corresponding threshold above.
The correspondence between conjuncts and thresholds is as follows.  The
rounding, \(0<m\), \(0\le p\le1/2\), \(K\le n\), and \(1\le2pm\) assertions
are supplied by \(N_{\mathrm{pos}}\), (R), and (P).  The small-pair, fiber,
expectation-absorption, McDiarmid-union, terminal-small-pair, logarithmic, and
huge-threshold size conjunct \(t_1k\le\varepsilon_1k^2\) are supplied
respectively by \(N_1,N_2,N_3,N_4,N_5,N_6,N_7\).  The huge-overlap count
\[
\binom{2k/t_1+1}{2}_{\mathbb R}100L^3\le\frac12k
\]
is supplied by \(N_{7a}\).  The same-color huge size and projected-pair
estimates are supplied by \(N_{7b}\) and \(N_{7e}\); the projected-overlap
size and multiplicity corrections by \(N_{7c}\) and \(N_{7d}\).  The four
relative-error comparisons with \(\varepsilon_2/100\), namely those with
left sides \(2pm\), \((2k/t_1+1)100L^3\), \((2k/t_1+1)200L^3\), and
\((2k/t_1+1)400L^5\), are supplied by \(N_{7f},N_{7g},N_{7h}\), while the two
pure \(\varepsilon_2\)-conjuncts were built into the initial choice of
\(\varepsilon_2\).  Finally,
\[
N_8,N_9,N_{9a},N_{10},N_{11},N_{12},N_{13},N_{14},N_{15},N_{16}
\]
supply, in order, the large-cutoff scale estimate, the large-versus-huge
cutoff, the two-term cutoff comparison with \((1+\varepsilon_2)\), the
large-cutoff lower bound, the large-huge count comparison, the \(t_3^2k\)
estimate, the projected-overlap loss estimate, the combined coefficient
inequality, and the two exponential comparisons.  This lists every displayed
conjunct in the current
\(\noderef{ParameterHierarchyEventualComparisons}\) package.

It remains to assemble the fixed hierarchy inequalities.  The choices
give
\[
0<\varepsilon_2<\varepsilon_1<\eta=\varepsilon'<1,\qquad
\varepsilon'<\frac1{16},
\]
\[
3\varepsilon_2\le\frac{\varepsilon_1}{10},\qquad
8\sqrt{\varepsilon_1}+10\varepsilon_1\le\frac{(\varepsilon')^3}{2}.
\]
Together with \(n_0\ge\varepsilon_2^{-2}\) and the just-proved eventual
comparisons, this is exactly
\[
\noderef{ParameterHierarchy}(\varepsilon',\varepsilon_1,\varepsilon_2,n_0).
\]
The initial choice of \(\varepsilon'\) also gives
\[
0<\varepsilon',\qquad \varepsilon'<\varepsilon,\qquad \varepsilon'\le\varepsilon,
\]
so the existential conclusion follows.
\end{proof}
